% from ega0-1-fr.tex
\oldpage[0\textsubscript{I}]{58}

\begin{env}[6.4.2]
\label{0.6.4.2-fr}
Un module libre $\neq\{0\}$ est fidèlement plat~; il en est de même de la
somme directe d'un module plat et d'un module fidèlement plat. Si $S$ est
une partie multiplicative de $A$, $S^{-1}A$ n'est un $A$-module fidèlement
plat que si $S$ est formé d'éléments inversibles (donc $S^{-1}A=A$).
\end{env}

\begin{env}[6.4.3]
\label{0.6.4.3-fr}
Soit
\[
  0\longrightarrow M'\longrightarrow M\longrightarrow M''\longrightarrow 0
\]
une suite exacte de $A$-modules~; si $M'$ et $M''$ sont plats, et si l'un
d'eux est fidèlement plat, alors $M$ est fidèlement plat.
\end{env}

\begin{env}[6.4.4]
\label{0.6.4.4-fr}
Soient $A$ et $B$ deux anneaux, $M$ un $A$-module, $N$ un
$(A,B)$-bimodule. Si $M$ est fidèlement plat et si $N$ est un $B$-module
fidèlement plat, alors $M\otimes_A N$ est un $B$-module fidèlement plat. En
particulier, si $M$ et $N$ sont deux $A$-modules fidèlement plats, il en est
de même de $M\otimes_A N$. Si $B$ est une $A$-algèbre et si $M$ est un
$A$-module fidèlement plat, le $B$-module $M_{(B)}$ est fidèlement plat.
\end{env}

\begin{env}[6.4.5]
\label{0.6.4.5-fr}
Si $M$ est un $A$-module fidèlement plat et si $S$ est une partie
multiplicative de $A$, $S^{-1}M$ est un $S^{-1}A$-module fidèlement plat,
puisque $S^{-1}M=M\otimes_A(S^{-1}A)$ (6.4.4). Inversement, si pour tout
idéal maximal $\mathfrak{m}$ de $A$, $M_\mathfrak{m}$ est un
$A_\mathfrak{m}$-module fidèlement plat, $M$ est un $A$-module fidèlement
plat, car $M$ est $A$-plat (6.3.3), et on a
\[
  M_\mathfrak{m}/\mathfrak{m}M_\mathfrak{m}
  =(M\otimes_A A_\mathfrak{m})
  \otimes_{A_\mathfrak{m}}
  (A_\mathfrak{m}/\mathfrak{m}A_\mathfrak{m})
  =M\otimes_A(A/\mathfrak{m})
  =M/\mathfrak{m}M,
\]
donc l'hypothèse entraîne que $M/\mathfrak{m}M\neq0$ pour tout idéal
maximal $\mathfrak{m}$ de $A$, ce qui prouve notre assertion (6.4.1).
\end{env}

\subsection{Restriction des scalaires}
\label{subsection:0.6.5-fr}

\begin{env}[6.5.1]
\label{0.6.5.1-fr}
Soient $A$ un anneau, $\varphi:A\to B$ un homomorphisme d'anneaux faisant
de $B$ une $A$-algèbre. Supposons qu'il existe un $B$-module $N$ qui soit un
$A$-module \emph{fidèlement plat}. Alors, pour tout $A$-module $M$,
l'homomorphisme $x\mapsto 1\otimes x$ de $M$ dans
$B\otimes_A M=M_{(B)}$ est \emph{injectif}. En particulier, $\varphi$ est
injectif~; pour tout idéal $\mathfrak{a}$ de $A$, on a
$\varphi^{-1}(\mathfrak{a}B)=\mathfrak{a}$~; pour tout idéal maximal (resp.
premier) $\mathfrak{m}$ de $A$, il existe un idéal maximal (resp. premier)
$\mathfrak{n}$ de $B$ tel que $\varphi^{-1}(\mathfrak{n})=\mathfrak{m}$.
\end{env}

\begin{env}[6.5.2]
\label{0.6.5.2-fr}
Lorsque la condition de (6.5.1) est remplie, on identifie $A$ à un
sous-anneau de $B$ par $\varphi$ et plus généralement, pour tout $A$-module
$M$, on identifie $M$ à un sous-$A$-module de $M_{(B)}$. On notera que si
$B$ est alors \emph{noethérien}, il en est de même de $A$, car l'application
$\mathfrak{a}\mapsto\mathfrak{a}B$ est une injection croissante de
l'ensemble des idéaux de $A$ dans l'ensemble des idéaux de $B$~;
l'existence d'une suite infinie strictement croissante d'idéaux de $A$
entraînerait donc l'existence d'une suite analogue d'idéaux de $B$.
\end{env}

\subsection{Anneaux fidèlement plats}
\label{subsection:0.6.6-fr}

\begin{env}[6.6.1]
\label{0.6.6.1-fr}
Soit $\varphi:A\to B$ un homomorphisme d'anneaux faisant de $B$ une
$A$-algèbre. Les cinq propriétés suivantes sont équivalentes~:
\begin{enumerate}
  \item[$a)$] $B$ est un $A$-module fidèlement plat (autrement dit,
    $M_{(B)}$ est un foncteur exact et fidèle en $M$).
  \item[$b)$] L'homomorphisme $\varphi$ est injectif et le $A$-module
    $B/\varphi(A)$ est plat.

